% Generated by generate.py: do not edit.
\chapter{Part II: The Graduate Toolbox}
\section{Tome 6: Spaces}
\subsection*{1.\ Sets, Choice and the Real Numbers}
\begin{results}
\item[1.2] Theorem (Peano's axioms)
\item[1.3] Theorem (recursion)
\item[1.11] Theorem (uniqueness of $\R$)
\item[1.14] Theorem (Cantor, 1891)
\item[1.15] Theorem (Schr\"oder--Bernstein)
\item[1.19] Theorem (Zorn's lemma)
\end{results}
\subsection*{2.\ Metric Spaces and Completeness}
\begin{results}
\item[2.8] Lemma (extension)
\item[2.9] Theorem (completion)
\item[2.11] Theorem (Banach, 1922)
\item[2.13] Lemma (Cauchy product)
\item[2.19] Theorem (Hutchinson, 1981)
\end{results}
\subsection*{3.\ Topological Spaces}
\begin{results}
\item[3.7] Proposition (universal property of quotients)
\item[3.12] Corollary (intermediate value theorem)
\item[3.16] Proposition (Urysohn functions in metric spaces)
\end{results}
\subsection*{4.\ Compactness}
\begin{results}
\item[4.4] Theorem (Heine--Borel)
\item[4.5] Corollary (extreme values)
\item[4.9] Theorem (Heine)
\item[4.11] Theorem (Tychonoff, 1930)
\item[4.15] Theorem (Arzel\`a--Ascoli)
\item[4.19] Theorem (Stone--Weierstrass)
\end{results}
\subsection*{5.\ Baire's Theorem}
\begin{results}
\item[5.2] Theorem (Baire, 1899)
\item[5.7] Theorem (Banach, Mazurkiewicz, 1931)
\item[5.9] Theorem (Baire, 1899)
\end{results}
\subsection*{6.\ Measures}
\begin{results}
\item[6.4] Theorem (Dynkin's $\pi$-$\lambda$ theorem)
\item[6.5] Theorem (uniqueness of measures)
\item[6.7] Theorem (Carath\'eodory, 1914)
\item[6.10] Theorem (Carath\'eodory's extension theorem)
\item[6.11] Theorem (Lebesgue measure)
\item[6.12] Proposition (regularity)
\item[6.14] Proposition (quantile transform)
\item[6.15] Theorem (Vitali, 1905)
\item[6.19] Theorem (Borel, 1909)
\end{results}
\subsection*{7.\ The Lebesgue Integral}
\begin{results}
\item[7.4] Theorem (monotone convergence)
\item[7.7] Theorem (Fatou's lemma)
\item[7.9] Theorem (dominated convergence)
\item[7.10] Theorem (the integral of a continuous function)
\item[7.11] Theorem (Lebesgue, 1902)
\item[7.16] Theorem (Vitali's convergence theorem)
\end{results}
\subsection*{8.\ Product Measures and Change of Variables}
\begin{results}
\item[8.4] Theorem (product measure)
\item[8.5] Theorem (Tonelli)
\item[8.6] Theorem (Fubini)
\item[8.8] Theorem (change of variables)
\item[8.10] Theorem (polar coordinates)
\item[8.11] Theorem (the $n$-ball)
\item[8.12] Theorem (Gaussian densities)
\end{results}
\subsection*{9.\ \texorpdfstring{$L^p$}{Lp} Spaces}
\begin{results}
\item[9.1] Lemma (Young's inequality)
\item[9.2] Theorem (H\"older's inequality)
\item[9.3] Theorem (Minkowski's inequality)
\item[9.4] Theorem (Riesz--Fischer)
\item[9.5] Theorem (Jensen's inequality)
\item[9.8] Proposition (continuity of translation)
\item[9.10] Theorem (approximate identities)
\end{results}
\subsection*{10.\ Differentiation and Radon--Nikodym}
\begin{results}
\item[10.2] Theorem (Hahn decomposition)
\item[10.3] Corollary (Jordan decomposition)
\item[10.4] Theorem (Radon--Nikodym)
\item[10.5] Theorem (Lebesgue decomposition)
\item[10.6] Corollary (the three parts of a law)
\item[10.7] Lemma (Vitali's covering lemma)
\item[10.8] Theorem (maximal inequality)
\item[10.9] Theorem (Lebesgue's differentiation theorem)
\item[10.11] Theorem (differentiation of measures)
\item[10.12] Theorem (Lebesgue, 1904)
\item[10.14] Theorem (fundamental theorem of calculus)
\end{results}
\subsection*{11.\ Measures on Metric Spaces}
\begin{results}
\item[11.1] Theorem (Ulam)
\item[11.2] Lemma (partitions of unity)
\item[11.3] Theorem (Riesz representation)
\item[11.4] Theorem (portmanteau)
\item[11.7] Theorem (Prokhorov)
\item[11.8] Theorem (Kolmogorov, 1933)
\item[11.10] Lemma (mass distribution principle)
\end{results}
\section{Tome 7: Analysis in Depth}
\subsection*{1.\ Banach Spaces}
\begin{results}
\item[1.2] Lemma (Riesz's lemma)
\item[1.3] Theorem (Riesz)
\item[1.4] Theorem (Hahn--Banach, analytic form)
\item[1.6] Theorem (Hahn--Banach, geometric form)
\item[1.7] Theorem (uniform boundedness, Banach--Steinhaus)
\item[1.9] Theorem (open mapping)
\item[1.10] Theorem (closed graph)
\item[1.11] Theorem (duality of $L^p$)
\item[1.13] Theorem (Riesz--Markov)
\item[1.16] Theorem (Banach--Alaoglu)
\end{results}
\subsection*{2.\ Hilbert Spaces}
\begin{results}
\item[2.1] Theorem (projection on a closed convex set)
\item[2.3] Theorem (Riesz--Fr\'echet)
\item[2.5] Theorem (orthonormal bases)
\item[2.9] Theorem (Kronecker--Weyl)
\item[2.10] Theorem (von Neumann's mean ergodic theorem)
\item[2.12] Theorem (Moore--Aronszajn)
\end{results}
\subsection*{3.\ Operators, Spectra and Banach Algebras}
\begin{results}
\item[3.4] Theorem (Fredholm alternative)
\item[3.5] Lemma (Neumann series)
\item[3.6] Theorem (Gelfand)
\item[3.9] Theorem (Gelfand, 1941)
\item[3.11] Theorem (Wiener's lemma, 1932)
\end{results}
\subsection*{4.\ The Spectral Theorem}
\begin{results}
\item[4.4] Theorem (spectral theorem, Hilbert 1906)
\item[4.6] Theorem (Courant--Fischer)
\item[4.9] Lemma (Green's function)
\item[4.11] Theorem (Sturm--Liouville)
\item[4.12] Corollary (Weyl's law on an interval)
\item[4.14] Theorem (Mercer, 1909)
\item[4.16] Theorem (continuous functional calculus)
\item[4.17] Corollary (square roots)
\end{results}
\subsection*{5.\ The Fourier Transform}
\begin{results}
\item[5.2] Theorem (Fourier inversion)
\item[5.4] Theorem (Plancherel, 1910)
\item[5.5] Theorem (Hermite functions)
\item[5.6] Lemma (three lines lemma)
\item[5.7] Theorem (Riesz--Thorin)
\item[5.8] Corollary (Hausdorff--Young)
\item[5.12] Theorem (sampling in Hilbert space)
\item[5.14] Theorem (Poisson summation)
\item[5.15] Corollary (theta of a lattice)
\end{results}
\subsection*{6.\ Distributions and Sobolev Spaces}
\begin{results}
\item[6.11] Theorem (Newton's potential)
\item[6.13] Theorem (Sobolev embedding)
\item[6.15] Theorem (Rellich)
\item[6.17] Theorem (Lax--Milgram, 1954)
\item[6.19] Proposition (C\'ea's lemma)
\end{results}
\subsection*{7.\ Fourier Analysis on Finite Abelian Groups}
\begin{results}
\item[7.2] Theorem (duality)
\item[7.6] Theorem (Poisson summation)
\item[7.7] Lemma (splitting along a subgroup)
\item[7.8] Theorem (fast Fourier transform)
\item[7.9] Theorem (Donoho and Stark, 1989)
\item[7.10] Lemma (Chebotarev, 1926)
\item[7.11] Theorem (Tao, 2005)
\item[7.14] Theorem (quadratic reciprocity)
\item[7.16] Theorem (Meshulam, 1995)
\item[7.17] Theorem (Blum, Luby and Rubinfeld, 1990)
\end{results}
\subsection*{8.\ Haar Measure, Duality and the \texorpdfstring{$p$}{p}-adic Numbers}
\begin{results}
\item[8.3] Theorem (uniqueness)
\item[8.4] Theorem (Markov and Kakutani, 1936--38)
\item[8.7] Theorem (Banach, 1923)
\item[8.9] Lemma (Hausdorff, 1914)
\item[8.10] Theorem (Hausdorff, 1914)
\item[8.11] Theorem (some dual groups)
\item[8.14] Theorem ($\Q_p$ is self-dual)
\item[8.15] Theorem (Fourier analysis on $\Q_p$)
\item[8.20] Theorem (adelic Poisson summation)
\end{results}
\subsection*{9.\ Curves, Growth and Conformal Maps}
\begin{results}
\item[9.1] Lemma (continuous argument)
\item[9.4] Theorem (no retraction; Brouwer)
\item[9.5] Lemma (crossing paths)
\item[9.6] Lemma (arcs do not separate)
\item[9.7] Theorem (Jordan curve theorem)
\item[9.8] Theorem (Cauchy's theorem, homology form)
\item[9.9] Corollary (residue theorem, general form)
\item[9.12] Theorem (Phragm\'en--Lindel\"of in a half-strip)
\item[9.13] Corollary (convexity)
\item[9.15] Theorem (the convexity bound)
\item[9.18] Theorem (Hadamard factorisation)
\item[9.19] Theorem (Montel, 1907)
\item[9.20] Lemma (injective holomorphic maps)
\item[9.21] Lemma (Schwarz; disc automorphisms)
\item[9.22] Theorem (Riemann mapping theorem)
\item[9.24] Theorem (Mellin inversion)
\item[9.26] Theorem (the maximum of geometric levels)
\end{results}
\subsection*{10.\ Tauberian Theorems}
\begin{results}
\item[10.1] Lemma (regular weights)
\item[10.2] Theorem (Abel, 1826)
\item[10.5] Theorem (Tauber, 1897)
\item[10.7] Theorem (Hardy, Littlewood, Karamata)
\item[10.8] Corollary (power series)
\item[10.9] Corollary (a converse to Abel)
\item[10.10] Lemma (local inverse)
\item[10.11] Theorem (Wiener, 1932)
\item[10.12] Theorem (Wiener's Tauberian theorem)
\item[10.13] Theorem (Wiener--Ikehara)
\item[10.14] Corollary (the prime number theorem)
\item[10.15] Corollary (progressions)
\end{results}
\section{Tome 8: Shapes}
\subsection*{1.\ Smooth Manifolds}
\begin{results}
\item[1.2] Proposition (chain rule)
\item[1.4] Lemma (mean value inequality)
\item[1.7] Theorem (inverse function theorem)
\item[1.8] Theorem (implicit function theorem)
\item[1.10] Theorem (four descriptions)
\item[1.21] Theorem (partitions of unity)
\item[1.25] Lemma (Hadamard)
\item[1.30] Theorem (regular value theorem)
\item[1.33] Theorem (Sard, equal dimensions)
\item[1.35] Theorem (Whitney, easy version)
\end{results}
\subsection*{2.\ Vector Fields and Flows}
\begin{results}
\item[2.1] Lemma (Gr\"onwall)
\item[2.2] Theorem (Picard--Lindel\"of)
\item[2.3] Proposition (escape lemma)
\item[2.4] Theorem (continuous dependence)
\item[2.5] Theorem (smooth dependence)
\item[2.10] Theorem (the bracket as a small loop)
\item[2.11] Theorem (Lie derivative)
\item[2.12] Theorem (commuting flows)
\item[2.14] Theorem (Frobenius)
\item[2.16] Theorem (hairy ball)
\end{results}
\subsection*{3.\ Differential Forms and Stokes' Theorem}
\begin{results}
\item[3.5] Theorem (exterior derivative)
\item[3.12] Proposition (boundary orientation)
\item[3.13] Theorem (Stokes)
\item[3.15] Theorem (no smooth retraction)
\item[3.16] Theorem (Brouwer)
\item[3.18] Theorem (homotopy formula)
\item[3.19] Corollary (Poincar\'e lemma)
\item[3.21] Corollary (the winding number as an integral)
\end{results}
\subsection*{4.\ Curvature}
\begin{results}
\item[4.3] Proposition (Cartan's first structure equations)
\item[4.5] Proposition (orthogonal coordinates)
\item[4.9] Theorem (Gauss, Theorema Egregium)
\item[4.11] Theorem (Gauss--Bonnet--Chern)
\item[4.13] Corollary (Gauss--Bonnet)
\item[4.18] Proposition (Green's formula)
\item[4.19] Theorem (the spectrum of the square flat torus)
\end{results}
\subsection*{5.\ The Fundamental Group and Coverings}
\begin{results}
\item[5.7] Lemma (uniqueness of lifts)
\item[5.8] Theorem (homotopy lifting)
\item[5.12] Theorem (lifting criterion)
\item[5.16] Theorem (classification of coverings)
\item[5.18] Theorem (Seifert--van Kampen)
\item[5.24] Theorem (Nielsen--Schreier)
\end{results}
\subsection*{6.\ Homology}
\begin{results}
\item[6.2] Lemma (the zigzag lemma)
\item[6.3] Lemma (the five lemma)
\item[6.7] Theorem (homotopy invariance)
\item[6.8] Theorem (small chains)
\item[6.9] Theorem (Mayer--Vietoris)
\item[6.10] Corollary (excision)
\item[6.18] Theorem (invariance of domain; Brouwer, 1912)
\item[6.20] Theorem (Borsuk--Ulam)
\item[6.21] Corollary (the ham sandwich theorem)
\end{results}
\subsection*{7.\ Matrix Lie Groups}
\begin{results}
\item[7.5] Theorem (von Neumann, Cartan)
\item[7.8] Theorem (the spin double cover)
\item[7.9] Theorem (maximal torus and Weyl group of $U(n)$)
\end{results}
\subsection*{8.\ Haar Measure and Weyl Integration}
\begin{results}
\item[8.1] Theorem (Haar measure of a matrix Lie group)
\item[8.3] Theorem (uniqueness)
\item[8.7] Theorem (Gaussian sampling of Haar measure)
\item[8.8] Theorem (Weyl, 1925)
\item[8.10] Lemma (the Jacobian)
\item[8.12] Lemma (Andreief, 1883)
\item[8.13] Theorem (Heine's identity)
\item[8.14] Lemma (Gaudin)
\item[8.15] Theorem (correlation functions of random unitary matrices)
\item[8.16] Corollary (Dyson's pair correlation)
\end{results}
\subsection*{9.\ The Hyperbolic Plane and the Modular Group}
\begin{results}
\item[9.2] Theorem (the isometry group)
\item[9.3] Proposition (geodesic polar coordinates)
\item[9.4] Theorem (Iwasawa decomposition)
\item[9.5] Theorem (the geodesic flow is a matrix product)
\item[9.6] Theorem (elliptic, parabolic, hyperbolic)
\item[9.7] Proposition (the Farey tessellation)
\item[9.8] Theorem (Series, 1985)
\end{results}
\subsection*{10.\ Integral Geometry}
\begin{results}
\item[10.1] Theorem (uniqueness of the measure on lines)
\item[10.2] Corollary (Bertrand's chord, answered)
\item[10.3] Lemma (counting preimages)
\item[10.4] Theorem (Crofton, 1868)
\item[10.5] Corollary (Buffon's noodle)
\item[10.7] Theorem (Crofton and Cauchy)
\item[10.8] Theorem (Poincaré, 1896)
\item[10.9] Theorem (Cauchy's surface area formula)
\end{results}
\section{Tome 9: Structures II}
\subsection*{1.\ Categories}
\begin{results}
\item[1.9] Theorem (the double dual is natural)
\item[1.10] Theorem ($V$ is not naturally its dual)
\item[1.15] Proposition (universal property of abelianisation)
\item[1.16] Theorem (Yoneda's lemma)
\item[1.17] Corollary (an object is known by its maps)
\item[1.22] Theorem (uniqueness of adjoints)
\item[1.23] Theorem (right adjoints preserve products)
\end{results}
\subsection*{2.\ Tensors and Exactness}
\begin{results}
\item[2.3] Proposition (splitting lemma)
\item[2.4] Theorem (snake lemma)
\item[2.5] Corollary (short five lemma)
\item[2.6] Proposition (Hom is left exact)
\item[2.12] Theorem (tensor and Hom are adjoint)
\item[2.13] Corollary (tensor is right exact)
\item[2.16] Theorem (Baer's criterion)
\end{results}
\subsection*{3.\ Noetherian Rings and Hilbert's Theorems}
\begin{results}
\item[3.3] Theorem (Hilbert's basis theorem, 1890)
\item[3.5] Proposition (the Zariski topology)
\item[3.8] Lemma (Artin and Tate)
\item[3.9] Theorem (Zariski's lemma)
\item[3.10] Theorem (weak Nullstellensatz)
\item[3.11] Theorem (Hilbert's Nullstellensatz, 1893)
\item[3.12] Corollary (the dictionary)
\item[3.14] Lemma (Dickson)
\item[3.15] Proposition (division)
\item[3.19] Theorem (Buchberger's criterion, 1965)
\item[3.20] Corollary (Buchberger's algorithm)
\end{results}
\subsection*{4.\ Localisation and Dimension}
\begin{results}
\item[4.2] Proposition (universal property)
\item[4.3] Proposition (localisation is exact)
\item[4.4] Theorem (primes of a ring of fractions)
\item[4.6] Theorem (Nakayama's lemma)
\item[4.10] Theorem (the primes of $\Z\lbrack x\rbrack$)
\item[4.15] Theorem (Cohen and Seidenberg, 1946)
\item[4.18] Lemma (Noether's change of variables)
\item[4.19] Theorem (Noether normalisation, 1926)
\end{results}
\subsection*{5.\ Dedekind Domains}
\begin{results}
\item[5.6] Theorem (Noether, 1927)
\item[5.7] Corollary (to contain is to divide)
\item[5.11] Proposition (one and a half generators)
\item[5.14] Theorem (the fundamental identity)
\end{results}
\subsection*{6.\ Homological Algebra}
\begin{results}
\item[6.2] Theorem (comparison theorem)
\item[6.3] Lemma (horseshoe lemma)
\item[6.6] Theorem (long exact sequences)
\item[6.12] Theorem (universal coefficient theorems)
\end{results}
\subsection*{7.\ Group Cohomology}
\begin{results}
\item[7.4] Theorem (Hilbert, Noether, Speiser)
\item[7.5] Corollary (Theorem~90 of the \emph{Zahlbericht})
\item[7.7] Theorem (Schreier)
\item[7.11] Theorem (Schur, 1904; Zassenhaus, 1937)
\end{results}
\subsection*{8.\ Semisimple Algebras and Induction}
\begin{results}
\item[8.3] Theorem (Wedderburn, 1908)
\item[8.7] Theorem (Frobenius reciprocity, 1898)
\item[8.8] Proposition (the induced character)
\item[8.10] Theorem (Mackey, 1951)
\item[8.11] Corollary (Mackey's irreducibility criterion)
\item[8.13] Lemma (Burnside)
\item[8.15] Theorem (Burnside, 1904)
\item[8.16] Theorem (Frobenius, 1901)
\end{results}
\subsection*{9.\ The Symmetric Group}
\begin{results}
\item[9.4] Lemma (Dominance lemma)
\item[9.6] Theorem (James' submodule theorem, 1978)
\item[9.7] Theorem (Young, Frobenius, Specht)
\item[9.11] Theorem (Robinson--Schensted correspondence)
\item[9.12] Theorem (Schensted, 1961)
\item[9.14] Theorem (Hook length formula; Frame, Robinson and Thrall, 1954)
\item[9.15] Theorem (Frobenius, 1900)
\end{results}
\subsection*{10.\ Shuffling Cards}
\begin{results}
\item[10.2] Theorem (Fourier inversion and Plancherel)
\item[10.4] Lemma (Diaconis--Shahshahani upper bound lemma)
\item[10.8] Theorem (Diaconis and Shahshahani, 1981)
\item[10.9] Corollary (Cutoff for random transpositions)
\end{results}
\subsection*{11.\ Compact Groups}
\begin{results}
\item[11.2] Theorem (Schur orthogonality)
\item[11.3] Theorem (Peter--Weyl, 1927)
\item[11.6] Theorem (Weyl's character formula for $U(n)$, 1925)
\item[11.8] Lemma (Pieri's rule for one box)
\item[11.9] Theorem (Diaconis--Shahshahani 1994, Rains 1998)
\end{results}
